\chapter{Formula Reference}
\label{ch:formularef}

\section{Starting a formula}
\index{formulas!syntax}

A cell entry is a formula if, and only if, its \emph{first} character is
\lit{=}.

\begin{ccexamples}[]{0.44\linewidth}
\exn{=1+1}{2}{a formula}
\exn{\textvisiblespace=1+1}{=1+1}{a leading space, so it is text}
\exn{+A1}{+A1}{text --- Lotus notation is not supported}
\exn{@SUM(A1:A9)}{@SUM(A1:A9)}{text}
\end{ccexamples}

\section{Operators}
\index{operators}

There are twelve, in five levels of precedence.

\begin{ccfigure}
\begin{tikzpicture}[x=1pt,y=-1pt,
   op/.style={font=\ttfamily\fontsize{9.5}{11.5}\selectfont,text=ccink},
   dsc/.style={font=\sffamily\fontsize{8}{10}\selectfont,text=ccink,anchor=west},
   lv/.style={font=\sffamily\bfseries\fontsize{8}{10}\selectfont,text=ccink,
              anchor=east}]
  \foreach \i/\lv/\ops/\d in {%
     0/{binds tightest}/{$-$x~~+x~~(~)}/{unary sign, parentheses},
     1/{}/{\^{}}/{power, left to right},
     2/{}/{*~~/}/{multiply, divide},
     3/{}/{+~~$-$}/{add, subtract},
     4/{binds loosest}/{=~~<>~~<~~>~~<=~~>=}/{comparison, does not chain}}{
    \pgfmathsetmacro{\yy}{\i*26}
    \fill[black!6] (52,{\yy-9}) rectangle (330,{\yy+11});
    \node[lv] at (46,{\yy+4}) {\lv};
    \node[op,anchor=west] at (58,{\yy+4}) {\ops};
    \node[dsc] at (190,{\yy+4}) {\d};
  }
  \draw[ccink,line width=0.6pt,-{Latex[length=4pt]}] (338,-4) -- (338,114);
  \node[font=\sffamily\fontsize{7}{8.4}\selectfont,text=ccink,
        rotate=-90,anchor=south] at (344,55) {evaluated later};
\end{tikzpicture}
\caption{Operator precedence, tightest binding at the top}
\label{fig:precedence}
\end{ccfigure}

\subsection{Consequences worth knowing}

\begin{ccexamples}[]{0.44\linewidth}
\exn{=2+3*4}{14}{multiplication first}
\exn{=(2+3)*4}{20}{}
\exn{=100/10/2}{5}{left to right, not 100/(10/2)}
\exn{=2\^{}3\^{}2}{64}{power is \emph{left}-associative, as in Excel}
\exn{=-2\^{}2}{4}{the sign binds tighter than the power}
\exn{=1+1=2}{TRUE}{arithmetic before comparison}
\end{ccexamples}

\begin{ccnote}
\fml{=2\^{}3\^{}2} is 64 here and 64 in Excel, but 512 in ordinary mathematical
notation and in most programming languages. \fml{=-2\^{}2} is 4 here and 4 in
Excel, but $-4$ in mathematics. \ccprod{} agrees with the spreadsheet
convention, not the mathematical one, on purpose.
\end{ccnote}

\subsection{Comparison operators}

Written \lit{=}, \lit{<>}, \lit{<}, \lit{>}, \lit{<=}, \lit{>=}. There is no
\lit{!=} and no \lit{==}. They produce \msg{TRUE} or \msg{FALSE}.

Comparisons do not chain: \fml{=A1<B1<C1} is a syntax error and the entry is
refused.

\begin{cccaution}
\textbf{Comparing text does not work properly.} Two cells that both contain
\lit{abc} compare as \emph{not} equal, because \ccprod{} compares text by
identity rather than by content and each piece of text in the workbook is a
separate object. \fml{="a"="a"} inside a single formula is \msg{FALSE} for the
same reason.

Worse, unequal text compares as ``greater'' in both directions, so
\fml{="a">"b"} and \fml{="b">"a"} are both \msg{TRUE} and \lit{<} on text is
always \msg{FALSE}.

Comparison is only reliable on numbers, dates and booleans. Do not build a
worksheet around comparing text.
\end{cccaution}

\subsection{What is missing}

There is no text concatenation operator --- no \lit{\&}. There is no percent
postfix: \fml{=50\%} is a syntax error, not 0.5.

\section{Cell references}
\index{references!syntax}

A reference is one or two column letters followed by one to five row digits:
\cellref{A1}, \cellref{B9}, \cellref{AZ100}, \cellref{IV65535}. Letters may be
upper or lower case.

Columns run \cellref{A} to \cellref{IV} (1 to 256) and rows 1 to 65,535.
Anything outside that --- \fml{=A0}, \fml{=ZZ1}, \fml{=A70000} --- is an invalid
reference and the entry is refused.

\subsection{Dollar signs: absolute references}
\index{absolute references}\index{references!absolute}\index{\$ (dollar sign)}

A \lit{\$} before the column letters, before the row digits, or before both,
stops that part of the reference moving when the formula is copied.

\begin{cctable}{Written}{When the formula is copied}{1.4in}{3.42in}
\fml{B1} & Relative. Both the column and the row move. \\
\fml{\$B1} & The column stays at \cellref{B}; the row moves. \\
\fml{B\$1} & The row stays at 1; the column moves. \\
\fml{\$B\$1} & Absolute. Always \cellref{B1}, wherever it is copied. \\
\end{cctable}

This is what makes one formula worth copying down a column. Put
\fml{=B2*\$C\$1} in a cell and copy it down: each row multiplies its own
\cellref{B} by the single rate in \cellref{C1}, because the first reference
moves and the second does not.

\begin{ccnote}
The evaluator does not know the difference. The compiler parses \lit{\$} and
discards it, so \fml{=\$B\$1} and \fml{=B1} are the same bytecode and give the
same answer. The \lit{\$} survives in the formula's \emph{source text}, which
is what the rewriter reads when a cell moves --- and the rewriter is the only
thing that needs to know.

One consequence: \textbf{a \lit{\$} does not protect a reference from an
inserted or deleted row}. A dollar sign fixes a reference against copying, and
a row that has physically moved has moved for everybody. See
Chapter~\ref{ch:rows}.
\end{ccnote}

\subsection{Ranges}

Two references and a colon: \cellref{A1:B9}. No spaces are allowed around the
colon --- \fml{=A1 : B2} is a syntax error, although \fml{=A1 + B2} is fine.

A range is always a rectangle, not a list of cells. \fml{=SUM(B1:C2)} adds four
cells.

The corners may be given either way round: \cellref{B20:A1} and
\cellref{A1:B20} are the same rectangle.

\begin{cccaution}
\textbf{Only five functions understand a range}: \fn{SUM}, \fn{AVERAGE},
\fn{MIN}, \fn{MAX} and \fn{COUNT}. Give a range to any other function and you
get \errv{\#VALUE!} --- \fml{=ABS(A1:A3)} and \fml{=ROUND(A1:A3,2)} both fail
that way.

A range used outside a function collapses to its top-left cell without
comment: \fml{=A1:B9*2} quietly means \fml{=A1*2}.
\end{cccaution}

There are no whole-column references (\lit{A:A}), no three-dimensional ranges
across worksheets, and no named ranges.

\subsection{References to another worksheet}

Three spellings:

\begin{ccexamples}[]{0.44\linewidth}
\exn{=Q1!B5}{}{plain name}
\exn{=SUM(Q1!B2:B40)}{}{a range}
\exn{='My Data'!A1}{}{quotes for a name with a space}
\end{ccexamples}

An unquoted name may contain letters, digits, underscores and full stops.
Anything else --- including a space, or a leading digit --- needs single quotes.
Names are matched without regard to case and must match the whole stored name.

The prefix applies to \textbf{one following reference only}. In
\fml{='My Data'!A1+A3}, \cellref{A3} is on the current worksheet.

A worksheet name that does not exist makes the reference invalid and the entry
is refused.

\section{Values in a formula}

\subsection{Numbers}

Digits, an optional decimal point, and an optional exponent introduced by
\lit{e} or \lit{E} with an optional sign: \lit{42}, \lit{3.14}, \lit{1.5E2},
\lit{15e-1}. At most twenty-three characters.

\subsection{Text}

In double quotes. A doubled \lit{""} inside is one literal quote.

\begin{ccexamples}[]{0.44\linewidth}
\ex{="hello"}{hello}
\ex{="say ""yes"""}{say "yes"}
\end{ccexamples}

A string literal in a formula may be at most \textbf{forty-eight} characters.
Longer than that, or unterminated, is a syntax error.

\subsection{TRUE and FALSE}

Written without quotes, in any case.

\section{Spaces}

Only the space character is skipped, and only in certain places: around
operators, around parentheses, between a function name and its opening bracket,
and around commas. \fml{=~~1~~+~~2~~} is a valid formula.

Tabs are not skipped and neither are spaces inside a range's colon. Anything
left over after a complete expression is a syntax error, so \fml{=1 2} is
refused.

\section{Argument separator}

A comma. Semicolons are not accepted.

\section{How a cell is read by a formula}
\index{cells!coercion}

\begin{cctable}{The cell holds}{The formula sees}{1.75in}{3.07in}
A number or a date & That number \\
\msg{TRUE} or \msg{FALSE} & 1 or 0 \\
Text & Text. In arithmetic that is \errv{\#VALUE!}; inside a range it is
skipped \\
An error & The same error, which propagates \\
A formula & That formula's last calculated value \\
Nothing at all & \textbf{Zero} when read on its own --- \fml{=Z90+1} is 1 ---
but \textbf{skipped} when it falls inside a range \\
\end{cctable}

That last row is the one to remember. \fml{=AVERAGE(B1:B10)} over three
populated cells divides by three, not by ten, which is what makes it agree with
Excel.

\section{Limits}
\index{formulas!limits}

\begin{cctable}{Limit}{Value}{2.6in}{2.22in}
Longest formula that can be typed & 80 characters \\
Longest formula in a file & 255 characters \\
Longest string literal & 48 characters \\
Values live at once during evaluation & 10 \\
Compiled size of one formula & 320 bytes \\
Nesting depth & No fixed limit; bounded by the above \\
\end{cctable}

\begin{ccnote}
The ten-value evaluation limit is the one you can reach by accident. A function
call with eleven or more arguments exceeds it, and the cell shows
\errv{\#ERR11!} --- which is the program's way of saying it ran out of room,
not that your formula is wrong. Use a range instead of a long argument list:
\fml{=SUM(A1:A20)} costs one value where \fml{=SUM(A1,A2,A3,\ldots)} costs
twenty.
\end{ccnote}

\section{When a formula is refused}

These are found while the formula is being compiled --- as you press
\key{Return}. Nothing is stored and the cell keeps what it had. No message is
shown.

\begin{cctable}{Cause}{Example}{2.6in}{2.22in}
Syntax error & \fml{=1+}~~\fml{=(1}~~\fml{=1 2}~~\fml{=A1<B1<C1} \\
Unknown name or function & \fml{=NOSUCH(1)}~~\fml{=SUM} \\
Invalid cell reference & \fml{=A0}~~\fml{=ZZ1}~~\fml{=Nowhere!A1} \\
Wrong number of arguments & \fml{=IF(1)}~~\fml{=ABS(1,2)} \\
Too long & Over 255 characters, or over 320 bytes compiled \\
\end{cctable}

\section{When a formula produces an error}

These go into the cell and are displayed. Chapter~\ref{ch:messages} lists them all.

\begin{cctable}{Value}{Cause}{1.2in}{3.62in}
\errv{\#DIV/0!} & Division by zero; \fn{MOD} by zero; \fn{AVERAGE} of nothing \\
\errv{\#VALUE!} & Text where a number was needed, or a range given to a
function that does not take one \\
\errv{\#NUM!} & Overflow, or an argument out of range \\
\errv{\#CYCLE!} & Circular reference \\
\errv{\#ERR}\emph{n}\errv{!} & A fault in the program rather than in the
formula, reported with its own code. \errv{\#ERR11!} is the one you can
provoke: more than ten live values \\
\end{cctable}
